\subsection{全量实验的索引、版本与证据合同}
任何一张图的一个核数配置是一次独立的优化单元。令索引集合 $\mathcal I=\{1,\ldots,100\}\times\{1,\ldots,5\}$，则每问均应有 $|\mathcal I|=500$ 条终态记录。每条记录至少保存图标识 $i$、核数 $k$、方案 $p_{i,k}$、官方完成时间 $T_{i,k}$、官方额外搬运 $D_{i,k}$，以及用于重放的输入与评价器身份。只报告五个均值而不保存这 $500$ 条明细，无法检查平均公式、重复配置或隐匿负例。三问对应冻结证据版本 Q1--r02、Q2--r07、Q3--r03；逐例表直接由相应目录的 CSV/JSON 记录生成，完整目录路径见附录的结果引用。原始赛题输入和固定硬件参数不在论文中被重写为“新数据”，后续若修改搜索器，应产生新的版本证据，而不能覆盖冻结旧结果。

正式结果的核对由三个相互独立的层次组成。\emph{方案层}检查非原始 COPY 算子的唯一覆盖、子图编号及每核顺序、商图与顺序弧无环，原图算子与张量属性未被改写。\emph{执行层}重新调用赛题评价程序，验证每个 Task 的 COPY、Pipe 依赖、容量及场景规则有合法时间线，并将最终周期和搬运与记录逐项比较。\emph{统计层}从逐例结果独立聚合五点平均加速比、总周期和搬运总量，检查单核固定基准是否一致。三层不能互相代替：CSV 行数为 $500$ 只证明覆盖索引，不证明其中每个方案可运行；官方输出可以证明该场景的执行结果，但若论文把总周期比错误写成平均加速比，统计结论仍不成立。

给每条记录附上内容哈希 $h_i=H(G_i)$、硬件配置哈希 $h_{s,k}=H(s,k,\vartheta)$、方案规范指纹 $h_p=H(\operatorname{canon}(p))$ 和评价器版本 $h_E$，则缓存键应至少包含
\begin{equation}\label{eq:evaluation-cache-key}
K=(h_i,h_{s,k},h_p,h_E).
\end{equation}
若改变场景 A/B/C 或 L2 容量、带宽，却复用旧键的评价结果，可能把不同问题的周期混淆；若只用方案名而不含完整映射与顺序，也可能把同名不同方案合并。式\eqref{eq:evaluation-cache-key}描述可复核的缓存合同，不表示所有历史开发运行都已按此键存储。交付结果以独立官方重放为准，缓存只用于减少确定相同配置的重复计算。

\subsection{指标从逐例记录到论文表图的推导}
设 $\mathbf T_k=(T_{1,k},\ldots,T_{100,k})$ 为某问固定场景下的周期向量，$\mathbf T_1^0$ 为题目规定的整图启发式单核基准。赛题五点曲线按式\eqref{eq:mean-speedup-derivation}逐元素相除再求平均：
\begin{equation}\label{eq:validation-vector-mean}
\overline S_k=\frac1{100}\mathbf 1^\top
\bigl(\mathbf T_1^0\oslash\mathbf T_k\bigr),
\end{equation}
其中 $\oslash$ 为逐元素相除。总周期效率为 $C_k=\mathbf 1^\top\mathbf T_k$；不同核数的总周期比 $C_1^0/C_k$ 是以多核周期为权重的平均比值，见式\eqref{eq:weighted-speedup}。例如问题一五核的 $\overline S_5=3.665903$，而总周期比约为 $312\,479\,373/73\,170\,501=4.270565$；两者差别并不是程序错误，而是权重口径不同。若把后者绘成赛题规定的平均加速比，将夸大这一指标。问题二五核对应的均值为 $4.159974$，总周期比为 $4.751102$；问题三五核 L2 优化方案均值为 $4.240375$，但同核无 L2 到 L2 的逐图时间比均值仅为 $1.022567$，这些数也不能直接相互替换。

对逻辑搬运，按核数汇总为 $D_k^{\Sigma}=\sum_iD_{i,k}$。由于 $D_{i,k}$ 的物理口径并不等于实际 DDR 服务字节，不能用 $D_k^{\Sigma}/60$ 与 $C_k$ 的比例作为拥塞强度的准确估计；在问题三尤其不能从 L2 命中字节直接减去 $D_k^{\Sigma}$。可以报告相对变化 $\Delta D_k=D_k^{\Sigma}(\text{new})-D_k^{\Sigma}(\text{base})$，但是否接受新方案仍按每配置字典序决策。聚合后周期变快而 $\Delta D_k>0$ 并不矛盾，因为新增搬运可能没有落在关键路径上；反过来，搬运下降也不保证更快。三问的附录表以每配置为最细粒度保留这种不一致，从而允许独立读者挑选任何一张图复核。

\subsection{成对比较、胜负台账与负例定义}
公平比较两个算法版本 $A,B$ 时，应在同一集合 $\mathcal J\subseteq\mathcal I$ 上建立成对差值 $\delta_j=T_j^{A}-T_j^{B}$，并记录
\begin{equation}\label{eq:paired-counts}
N_+=\sum_{j\in\mathcal J}\mathbf1\{\delta_j>0\},\quad
N_0=\sum_{j\in\mathcal J}\mathbf1\{\delta_j=0\},\quad
N_-=\sum_{j\in\mathcal J}\mathbf1\{\delta_j<0\}.
\end{equation}
同时给出总周期差 $\sum_j\delta_j$ 和搬运差。只报告 $N_+$ 会掩盖少数大幅退步，只报告总差会掩盖多数配置毫无变化。对开发集上的参数选择，若先观察结果再改变规则，再用同一小集合报告“独立验证”，会把调参和验证混同；全量固定规则回放和开发集消融必须明确区分。本文最终百图成绩不是随机抽样推断，也没有给固定一百图构造没有意义的“总体置信区间”；结论限定于题目给定集合和参数。

问题二 r05 到 r07 的最终对照是一例：在 $400$ 组多核配置上，$N_+=13$、$N_0=387$、$N_-=0$，周期和从 $403\,398\,480$ 降至 $403\,371\,635$，净减少 $26\,845$ cycles，约为 $0.006655\%$；额外搬运增加 $1\,267\,154$ bytes，约为 $0.028884\%$。由此能说 r07 在这批配置中取得小幅、全量可复核的周期改善，且搬运次指标存在代价；不能把 $13/400$ 描述成普遍的大收益，也不能因为净改善大于零就声称新结构候选总能更快。问题三的 $68$ 组“新增搜索改善”则是\emph{有 L2 场景内}的成对比较，与四组“最终 L2 比无 L2 慢”的跨场景负例属于不同参照，必须分开登记。前者包含保底所以没有场景内退步，后者说明硬件场景改变未保证固定方案的单调性。

若需评价候选排序而不是最终算法，应在\emph{同一候选池}中保存所有 $\widehat T_j,T_j$。Spearman 等级相关系数可用于度量总体单调关系，但在正式预算 $B$ 很小的场景中，更直接的是前两名覆盖率与式\eqref{eq:selection-regret}的选择遗憾。定义最优候选被前二名覆盖的比例为 $100^{-1}\sum_i\mathbf1\{j_i^*\in\{\widehat j_{i,1},\widehat j_{i,2}\}\}$，前提是每配置确实有至少两个\emph{不同}的、均已官方评价的候选。若很多配置只产生一个独特候选，不能把它当成排序器的完美命中。本文没有在缺完整候选池官方真值的地方填造相关系数；已有代理预测改善但正式评价输赢接近随机的开发证据，只能证明需要继续校准，不是某个数值相关系数的证据。

\subsection{三问的约束回归与反事实检查}
问题一的核心回归是验证每个边界 Task 确实经 DDR 中转，即便两个 Task 被映射到同核也不保留私有缓存；场景 A 的 $100$ 周期间隔与跨核 $1000$ 周期释放分别落在正确的 Task 边。若误把场景 B 的同核驻留带入问题一，计算出的 $D$ 将系统性偏低。问题二的回归是验证同核所有子图组成一个 Task，同核消费者可在容量允许时复用，跨核输入在源端 	exttt{COPY\_OUT} 完成后再延迟 $500$ 周期释放；若把 $500$ 加到整个 Task，周期会系统性偏高且错误阻塞不相关 Pipe。问题三再检查共享 L2 是全核单实例、只读、FIFO、发射时查找和完成时填充，并与 DDR 在不同带宽池服务；若把缓存命中当成已消除的逻辑 COPY，命中率与额外搬运的物理意义都会错位。

一个有力的反事实检查是\emph{固定方案改变单一场景机制}。问题二到问题三保存同一 $p^B$，先关闭 L2 得到 $T^B(p^B)$，再打开 L2 得到 $T^{L2}(p^B)$；之后才允许重新调度得 $T^{L2}(p^C)$。表\ref{tab:q3-total-cycle}的三列恰与这三步对应。问题一到问题二也可用式\eqref{eq:q2-scene-alg-decomp}设计类似对照，但两场景的 Task 构造差别大，须先验证移植方案在 B 中合法。若不做固定方案重放，就无法区分 L2 本身改善与搜索器改进；若只看最终 Q3 比 Q2 的平均加速比，则还可能误读两种统计权重。固定方案反事实不是现实硬件的受控实验，却是仿真模型中清楚的机制消融。

\subsection{核数非单调与结构解释的边界}
对同一图的相邻核数，定义 $\eta_{i,k}=T_{i,k}-T_{i,k-1}$。若 $\eta_{i,k}>0$，说明本次有限搜索在 $k$ 核配置给出的方案较 $k-1$ 核方案慢，而不是证明“多一个物理核心必然使最优时间更差”。通常可以让新核心闲置并复用旧的映射，但不同核数的固定评估配置、候选生成与预算若未包含嵌入旧方案，就不能直接从有限解得最优值单调性。问题二记录中的 $12$ 个相邻核数退步、问题三中的 $10$ 个，提示算法在某些图上存在映射或候选覆盖不足。示例包括第 $64$ 图问题二四核 $10\,160$、五核 $11\,391$ cycles，以及第 $71$ 图三核 $9\,867$、四核 $10\,732$ cycles。它们值得进一步分析，但不能只凭这两个数判定 DDR、同步或溢出中的哪一项负责；原因必须由具体时间线中关键输出的激活约束确认。

理论上若允许把旧 $k-1$ 核方案原样嵌入 $k$ 核且新核心闲置，二者使用相同场景参数和评价语义，则全局最优满足 $T^*_{i,k}\le T^*_{i,k-1}$。推导是 $k-1$ 核的每个合法方案都属于 $k$ 核可行集，因此后者最小值不会增大。这个集合包含关系\emph{不}自动保证交付的有界启发式输出也单调，因为其候选池 $\mathcal C_{i,k}$ 未必包含上一核数的嵌入方案。一个可行的后续增强是将已验证的 $k-1$ 核方案作为 $k$ 核保底，重新评价后允许其与本核新方案竞争；但在未对百图证明输入输出接口完全同构、并检验新增保底的机会成本之前，本文不把它写成交付版已实现机制。

\subsection{粒度与内存压力的双向敏感性}
子图数量 $q$ 与最终周期的关系不应预设为单峰或单调。令 $T(q;p)$ 表示保留其它条件的某候选族中、具有 $q$ 个子图的方案时间。把一个子图拆成两个，至少可能改变三项：\emph{并行度}上限变高，\emph{边界流量}及场景 A 的 Task 间隔变多，\emph{核内峰值驻留}可能下降。三项交互，故 $T(q+1)-T(q)$ 的符号由最终事件时间线决定。若只把 $q$ 当作连续超参数做大规模扫描，也会违反题目“避免暴力迭代”的求解效率目标。本文使用整分量、少数拓扑窗口和小区域边界修复，形成离散而可解释的几类尺度，再在相同正式预算下比较。

内存压力也不能直接由“张量总字节超过容量”判断。对一个 Task，若每个张量的生命周期区间两两不重叠，即便总字节数远大于 L1/UB，也可能只需有限容量轮流处理；若许多张量生命周期重叠，总量较小也可能在局部峰值触发换出。理想存活峰值 $P_{c,r}=\max_tL^{\rm ideal}_{c,r}(t)$ 与面积 $A_{c,r}$ 分别刻画瞬时压力和长期占用，实际是否发生溢出还依赖核内启发式选择的次序。若保存方案的溢出搬运为零，继续以“减少溢出”为这批方案的主要方向就缺证据；若某些未选候选溢出较大，则它们仍可作为该机制的反例。问题三第 $44$ 图的收益来自\emph{已经发生的}溢出重读命中，是 L2 的场景效应，不表示问题一、二应追求更高溢出以便缓存命中。

\subsection{求解耗时与评估周期的不同量纲}
设 $t^{\rm wall}_{i,k}$ 为求解器在 CPU 上从读取输入到输出方案的墙钟秒数，$T_{i,k}$ 为方案在模拟 NPU 上的执行周期。这两个量没有直接换算关系。增加官方评价次数可能降低 $T$，却提高 $t^{\rm wall}$；复用已经准备的核内子问题可能大幅降低 $t^{\rm wall}$，却使同一方案的 $T$ 完全不变。题目要求算法单用例尽量在五至十分钟给出高质量方案，因此两类指标都要报告，并说明计时区间：冷启动含输入解析、单核基准、候选生成、官方评价和最终复核；热搜索仅记录某一阶段。已记录的搜索阶段最长约 $364$ 秒，但这不包含全部冷启动与最终复核，不能据此宣称所有大图端到端低于十分钟。另一方面，大图的单核基准计算本身可达十余分钟，若把它与多核搜索混在一个数字里，也不能准确判断哪一部分最值得工程优化。

设一次完整官方评价平均耗时 $c_E$，候选准备 $c_P$，预算 $B$，则粗略墙钟成本为 $c_{\rm init}+B(c_P+c_E)+c_{\rm replay}$；真正的大图成本还随核内启发式和事件数变化。若某两个候选复用同一核内准备键，后者可减少 $c_P$ 而不改变 $c_E$；若重复方案通过式\eqref{eq:evaluation-cache-key}直接取得已校验的正式结果，连重复的 $c_E$ 也可避免。把节省的时间用于一个结构不同的候选，在保持正式评价机会数不变时可能提高结果质量，但是否如此仍须成对对照。本文将主机耗时控制作为独立的后续改进方向，而不把搜索运行更久所得的改善伪装为“算法本身更快”。

\subsection{适用范围、失败样例与可迁移性}
三个问题的结果建立在赛题给定的算子周期、张量字节数、L1/UB 容量、共享 DDR 带宽、场景同步规则及官方启发式调度上。若硬件改为不同 Pipe 组合，式\eqref{eq:universal-compute-bound}须按资源种类重写；若 DDR 改为非公平服务，式\eqref{eq:ddr-fluid-common}不再适用；若 L2 改为 LRU 或允许写入，式\eqref{eq:q3-fifo-next}和缓存状态推导须整体改写。即便硬件相同，换一批图的分量数、共享输入程度和关键链深度也可能改变结构候选的有效性。因此本方法可迁移的是“场景规则展开、结构初解、事件反馈、有界正式评价”的求解框架，而非从这 $100$ 图学出的某个固定分块粒度或收益百分比。

失败样例是检验模型诚实度的重要证据。问题一某些图在增核后出现周期退步；问题二存在 $12$ 组核数非单调；问题三仍有四组相对无 L2 更慢。这些现象并不否定所有启发式，只表明目前有限候选池和代理排序尚未穷尽结构空间。若下一轮扩展搜索，应先冻结现有候选池、辨别真正不同的方案和被新动作取代的基础候选，随后用相同官方预算对照其净收益；若直接扩大评价次数，结果改善无法归因于模型修正。论文最终交付的正面结论限定为：统一规则在给定百图和固定配置下输出经官方复核的合法方案，并获得所列的平均加速比、总周期和逻辑搬运；它不包含对全局最优、未见图泛化或真实芯片能耗的额外承诺。
